\section{機械化された性質と接続}

本節では，前節までの規則から得られる性質と，実行可能推論・動的安全性への接続をまとめる．
solver の個別補題，trace invariant，回帰一覧は \texttt{../docs/details.md} に分離する．

\subsection{SourceTyping 自身の性質}

\subsubsection*{Slot demand}
$S;\Omega\vdash\tau\preccurlyeq\rho\dashv S'$ が ordinary equality でないなら，resolved
expected type は必ず slot-headed である．
\[
S;\Omega\vdash\tau\preccurlyeq\rho\dashv S'
\Longrightarrow
(S;\Omega\vdash\tau\doteq\rho\dashv S')
\ \lor\
\exists\kappa,\upsilon.\;S\rho=\MatcherSlot{\kappa}{\upsilon}.
\]
特に $S\rho$ が matcher-headed なら ordinary alignment に限られる．これにより「coercion は
slot demand を観測した checking cut で一度だけ」が judgment の性質になる．

\subsubsection*{三状態の extension}
すべての raw demand-directed family は supply を単調に進め，出力 substitution を入力 substitution と局所
solve delta の chronological replay に因子化できる．対応する Origin family は raw derivation
と同じ構造を持ち，ledger transition が入力 ledger を保ち，fresh supply の範囲内だけを追加・
freeze することを保証する．公開型，pattern dual，bindings，hole ledger に現れる flexible
variable は終端 supply で有界である．

\subsubsection*{Solved form}
capability-only，target-only，paired，二段階 matcher／slot solve，one-way matchingの各局所操作は，
入力substitutionの冪等性を保存する．これをtraversal順に合成し，alignment familyと，expression，
expression list，checking，user／primitive／data pattern，pattern list，arm，clause，
pattern-constructor capabilityを含む全14 raw demand-directed familyについて
\[
S.\mathsf{Idempotent}\Longrightarrow S'.\mathsf{Idempotent}
\]
を証明済みである．公開判断は恒等substitutionから始まるため，rootの終端substitutionは無条件に
solved formである．

\subsubsection*{No guess と freeze}
exact MGU は constraint 外で恒等，constraint range 内，solved form である．さらに Origin
admissibility により rigid key を固定し，rename-only key の structural strengthening を拒否する．
constructor／fresh consumer の structural variable は局所 solve を許すが，選択的 export と
matcher finalization を越えた producer leaf は rename-only になる．$\mathtt{let}$ は終端
substitution 下の generalization stability も要求する．

\subsubsection*{回帰の現在地}
origin を追跡しない局所導出では，量化 producer capability を後続 solve が
$\Any$ へ強化する $\mathit{capFreezeProgram}$ と，let-generalized capability を強化する
$\mathit{letCapFreezeProgram}$ の問題が現れる．これらは現行 public
$\mathsf{SourceTyping}$ の正例ではない．該当する局所 delta が Origin-aware solve で拒否されることは
証明済みであり，プログラム全体についても
$\neg\mathsf{SourceTyping}(\mathit{capFreezeProgram})$ と
$\neg\mathsf{SourceTyping}(\mathit{letCapFreezeProgram})$ を閉じる negative regression を構成済みである．
or-pattern，delegating matcher，let-polymorphic producer の public positive Origin certificate も
構成済みである．

\subsection{実行可能推論からSourceTypingへ}

公開 inference は停止する raw traversal と有限の fail-closed validator の合成である．
\[
\mathsf{infer}(\Sighat,\Gamma,e)=
\mathsf{bind}(\mathsf{inferRaw}(\Sighat,\Gamma,e),
  \lambda r.\;\mathsf{if}\;\mathsf{wBridgeCheck}(\Sighat,r)
  \;\mathsf{then}\;\mathsf{some}\;r\;\mathsf{else}\;\mathsf{none}).
\]
raw traversal の executable ledger と demand-directed の intrinsic Origin certificate は同じ policy を別の
役割で記録する．前者は solver を実行時に fail closed にし，後者は関係的 derivation の solve
admissibility を証明する．validator は trace から binder-local instance，alignment，matcher
finalization，let generalization，freeze event を再検査する．成功したfuelled traversalは，
expression synthesis／checking，expression list，pattern，matcher，arm，clauseを含む10 familyの
相互帰納により，同じsupply，prevailing substitution，origin ledgerを持つraw demand-directed derivationへ
再構成される．append-only historyは各再帰呼出しを一つのroot終端へ接続し，validatorの証拠から
$\mathtt{let}$，matcher，pattern constructorのterminal auditを構成する．従って公開成功から
次を得る．
\[
\frac{\mathsf{infer}(\Sighat,\Gamma,e)=\mathsf{some}\;r}
{\mathsf{SourceTyping}(\Sighat,\Gamma,e,r.\mathsf{resolvedTarget})}.
\]
この定理は $\mathsf{TypingInvariant}$ を経由せず，caller-supplied bridge，history，
$\mathsf{InferenceInputWF}$を前提にしない．

これとは独立に，$S_r=r.\mathsf{state.prevailing}$として公開成功から
$\mathsf{TypingInvariant}(\Sighat,S_r\Gamma,e,r.\mathsf{resolvedTarget})$を再構成する既存経路もある．
$\mathsf{TypingInvariant}$ は source acceptance ではなく，solver stateを持たない内部invariantで
あり，動的メタ理論との直接接続に用いる．一般contextについて二つの経路を同一視する主張はしない．

\subsection{動的安全性}

freeze 済み signature の checker は，全schemeのbinder外にmetavariableがないことと動的整合性を
一つに束ねた$\mathsf{FrozenSigWF}$を確立する．function-valuedなexhaustiveness checkerだけは，
signature構築時に固定した実装との等式をchecker soundnessへ渡す．その下で，typed
environment における評価は $\mathsf{TypingInvariant}$ を $\mathsf{ValueTy}$ へ保存する．matching
machine について，typed state の一段保存，局所 progress，到達可能 state の保存，成功 branch
の substitution typing を証明済みである．空 successor 列は正当な match failure であり stuck
ではない．一般の program termination は主張しない．

\subsection{SourceTyping からの公開安全性}

source-facingな公開安全性定理は次の形で機械化済みである．
\[
\frac{
  \mathsf{FrozenSigWF}(\Sighat)
  \qquad
  \mathsf{SourceTyping}(\Sighat,\emptyset,e,\tau)
}{
  \mathsf{SafeResult}_{\mathsf{Source}}(\Sighat,e,\tau,\mathcal S_F)
}
\;\rn{Source-Safety}
\]
$\mathsf{SafeResult}_{\mathsf{Source}}$はLeanの$\mathsf{SourceTyping.SafeResult}$を表し，
$\mathsf{TypingInvariant}(\Sighat,\emptyset,e,\tau)$と
$\mathsf{CoreSafety}(\Sighat,\mathcal S_F)$を保持する．従って内部certificateを公開定理のpremiseに
露出せず，同じ公開型でevaluationとmatching machineの安全性を利用できる．

この定理のstate-erasure部分では，入力cutより前のorigin policyを保存する
supply-scopedな残余substitutionと，各局所出力を一つのroot終端substitutionへ結ぶ
$\mathsf{StateFactorization}$を使う．全14 Origin familyのfactorizationとconstructor-wise erasureを
一つの終端cutへ固定した相互帰納にまとめ，supply，prevailing substitution，ledgerを消去する．

局所cutではなく公開終端で再確認すべき追加情報はterminal auditに分離する．そのfactsは三種であり，
$\mathtt{let}$は終端context／value typeから再計算したgeneralizationの一致，matcherはterminal hole
capabilityから再計算したevidence，shape，clause capability，arm exhaustiveness，coverage，pattern
constructorは終端dual／capabilityの$\mathsf{CapCompatible}$を記録する．variable nodeは追加factsを
持たず，canonical $\mathsf{Scheme}$のfinite opening transportから終端lookupとの一致を直接導く．
相互erasureはこれらを使って
$\mathsf{TypingInvariant}$，pattern，arm，clause certificateを再構成する．matcher-to-slotは終端
$\mathsf{CapabilityDemand}$，slot-to-slotは終端型等式へ消去される．

従ってclosed signature上のclosed-programについて
\[
\Sighat.\mathsf{SchemesClosed}
\Longrightarrow
\mathsf{SourceTyping}(\Sighat,\emptyset,e,\tau)
\Longrightarrow
\mathsf{TypingInvariant}(\Sighat,\emptyset,e,\tau)
\]
が成立する．$\mathsf{TypingInvariant}$の存在をdemand-directed ruleのpremiseに置く循環的な定義は用いず，$\mathsf{SourceTyping}$ derivation自身の
solved-form保存，Origin history，terminal auditから射影する．このclosednessは低レベルerasureの
正確なpremiseとして残るが，公開callerは別に渡さない．$\mathsf{FrozenSigWF}$が
$\Sighat.\mathsf{SchemesClosed}$をfieldとして保持し，実行可能checkerが全tableについて検査する．

\subsection{受理完全性}

逆向の受理完全性は，$\mathsf{SourceTyping}$ derivationのOrigin treeとterminal auditを同時に再帰することで
証明する．demand-directedのexact solve witnessとexecutable solver resultを相互factorizationで対応させ，
fresh allocation，context normalization，producer protection，fuel boundを保ったfuelled traversalを，
expression，checking，pattern，arm，clauseの全familyに対して再構成する．

各局所runは成功等式とともにvalidator eventのcoverage extensionを返す．ordinary eventは
traversal自身から，$\mathtt{let}$ generalization，matcher finalization，pattern-constructor
compatibilityはterminal auditから得る．pattern constructorでは，demand-directed導出が記録するdual／capabilityと
実行trace中のoperandsを同一視せず，両者とそのbisimulationをpaired witnessとして保持する．rootで
これらを$\mathsf{PairedRootCertifiedSynthesis}$に束ね，pattern-constructor条件はpaired witnessから
直接，$\mathtt{let}$とmatcherの条件も各局所cutのdemand-directed／実行operandを結ぶpaired witnessから，
type／dual alignmentと合成して有限の$\mathsf{wBridgeCheck}$の全条件を導く．exact-state leafは
このchronologyの対角な特別場合として埋め込む．Originとauditは$\mathsf{Prop}$，concrete runは
$\mathsf{Type}$に属するため，fuelに対するstrong recursionは各cutのpaired runを
$\mathsf{Nonempty}$で返す．canonical initial cutで$\mathsf{PairedRootCertifiedSynthesis}$へ束ね，
公開facadeは受理命題の内部でだけその証明消去境界を開く．

公開定理はM4と同じglobal signature条件の下で
\[
\mathsf{SourceTyping}(\Sighat,\emptyset,e,\tau)
\Longrightarrow
\mathsf{FrozenSigWF}(\Sighat)
\Longrightarrow
\mathsf{infer}(\Sighat,\emptyset,e)\ne\mathsf{none}
\]
である．実際のLean定理$\mathtt{SourceTyping.infer\_isSome}$は一般のcontextに対して述べられる．
$\mathsf{FrozenSigWF}$からはterminal factsのbisimulation輸送に必要なscheme closednessとcanonical arm
checkerを取り出す．$\mathsf{RawSourceVisible}$，$\mathsf{FreezeCompatible}$，solver success，
validator bridgeなどの実装向け条件はcaller premiseに残らない．これはvalidator単体の無条件
完全性ではなく，terminal-audited demand-directed fragmentから生成したtraceに対する受理完全性である．
recursive list matcher，multiset matcher，pattern constructorを含む$\mathtt{matchAll}$の三例を，
caller-supplied paired rootを持たない公開回帰として機械検査している．

\subsection{受理同値とannotation-freeness}

executable soundnessと受理完全性を合成すると，一般contextについて
\[
\mathsf{FrozenSigWF}(\Sighat)
\Longrightarrow
\left(
\exists\tau.\,\mathsf{SourceTyping}(\Sighat,\Gamma,e,\tau)
\Longleftrightarrow
\mathsf{infer}(\Sighat,\Gamma,e)\ne\mathsf{none}
\right)
\]
を得る．Leanでは$\mathtt{Inference.sourceTypable\_iff\_infer\_isSome}$として機械化し，
$\Gamma=\emptyset$への特殊化を$\mathtt{Inference.annotation\_freeness}$として公開する．
sourceの$\mathsf{Expr}$にはtype-ascription constructorがないため，closed termに対して何らかの型で
$\mathsf{SourceTyping}$が成立するかを，型annotationを入力せず有限な公開推論で決定できる．対応する
$\mathtt{sourceTypableDecidable}$は$\mathsf{FrozenSigWF}$の証明の下で
$\mathsf{Decidable}(\exists\tau.\,\mathsf{SourceTyping})$を返す．

型を返す入口についても
\[
\mathsf{inferType}(\Sighat,\Gamma,e)=\mathsf{some}\;\tau
\Longrightarrow
\mathsf{SourceTyping}(\Sighat,\Gamma,e,\tau)
\]
を証明済みである．完全性方向では，任意に与えた$\mathsf{SourceTyping}$ derivationのtargetと返値型の構文的一致ではなく，
$\mathsf{inferType}$が返す何らかの型と，その返値自身の$\mathsf{SourceTyping}$ derivationを同時に得る．さらに同じ
$\Sighat,\Gamma,e$に対する二つの$\mathsf{SourceTyping}$ target $\tau_1,\tau_2$には，共通の決定的な実行target $\upsilon$と，
各$\tau_i$の全residual capability／target metavariable上で局所的に可逆な二sort変数renaming
$R_i$が存在し，$R_i\tau_i=\upsilon$となることを機械化済みである．

initial supply以前のmetavariableまで固定するより強い形は一般contextでは成立しない．例えば
$f:?0\to ?0, x:?1$の下の$f\,x$では，constraint $?1\doteq ?0$のexact MGUをどちら向きに取るかにより，
公開targetは$?0$にも$?1$にもなる．両導出はterminal auditまで含む回帰として機械化している．
従って一意性は全residual metaのrenamingを法として述べる．型のinstance preorder上のprincipalityは
別の結果であり，$\mathsf{TypingInvariant}$全体のprincipality反例をこのdemand-directed結果には流用しない．

\subsection{機械化の範囲}

全 expression／pattern／arm／clause の raw demand-directed family，それらに対応する intrinsic Origin family，
ledger の基本 transition，alignment の slot-demand，state replay，supply boundedness，全14 raw familyの
solved-form保存，terminal audit，terminal-fixedな$mathsf{TypingInvariant}$への相互射影，closed-programの
$\mathsf{SourceTyping}\Rightarrow\mathsf{TypingInvariant}$，公開 inferenceから$\mathsf{SourceTyping}$および
$\mathsf{TypingInvariant}$へのreconstruction，$\mathsf{SourceTyping.safe}$によるsource-facingな公開安全性，
concrete dynamic safety，demand-directed derivationから全実行traversal familyへの受理完全性，terminal auditから
validator event coverageへの輸送，公開の$\mathtt{SourceTyping.infer\_isSome}$，一般contextの受理同値，
closed annotation-freeness，$\mathtt{inferType}$返値soundness，source typabilityのdecidabilityは
Lean~4に機械化済みである．同じsourceの$\mathsf{SourceTyping}$ targetが全residual二sortmetaの局所renamingを法として
一意であることと，open contextでinitial metaを固定する強い形の反例も機械化済みである．
public freeze回帰も閉じている．
主定理と一般補題に $\mathtt{sorry}$，$\mathtt{admit}$，project-defined $\mathtt{axiom}$ はない．
